\documentclass{amsart}
\input{C:/Users/jzchr/MathDocuments/preamble_general.tex}

\title{132 HW \#3}
\author{Jalen Chrysos}

\begin{document}

\maketitle

\textbf{Problem 18}:\\

Using an equation obtained in chapter 22, the electric field distance $r$ away from a charged line of charge density $\lambda$ is $2k\lambda/r$, so
$$
840 = \frac{2k\lambda}{0.4} \implies \lambda = 1.87\cdot 10^{-8} C/m
$$
so in 2cm of wire there is
$$
1.87 \cdot 10^{-8} \cdot 0.02 = 3.7\cdot 10^{-10} C
$$
of charge.\\

\textbf{Problem 24}:\\

(a): Outside the solid sphere, electric field behaves just like a point charge $Q$, so 
$$
E = \frac{kQ}{r^2} \implies 940 = \frac{kQ}{0.08^2} \implies Q = 6.7 \cdot 10^{-10} C
$$
giving a charge density of 
$$
Q/V = \frac{6.7\cdot 10^{-10}}{\tfrac43\pi (0.04)^3} = 2.5 \cdot 10^{-6} C/m^3.
$$

(b): On the \textit{inside} of the sphere, the electric field is $kQr/R^3$, so
$$
E = \frac{kQr}{R^3} = \frac{kQ}{r^2} \cdot \Big(\frac{r}{R}\Big)^3 = 940 \cdot (1/2)^3 = 117.5 N/C.
$$

\textbf{Problem 40}:\\

(a): The electric field magnitude at radius $r$ is 
$$
E(r) = 
\begin{cases}
	kqr/a^3 & r<a\\
	kq/r^2 & a<r<b\\
	0 & b<r<c\\
	kq/r^2 & c<r
\end{cases}
$$
and it is pointed outward everywhere it is nonzero

(b): Here is the graph:
\begin{center}
	\includegraphics[width=8cm]{C:/Users/jzchr/MathDocuments/Class Work/Physics 132-133/132HW3/HW3_40.png}
\end{center}
Here I've set $a=4,b=8,c=12$.\\

(c): $-q$\\

(d): $q$\\

(e): The field lines are pointed outward at every point $r<b$ and $r>c$, with magnitude according to the graph above.
\\

\textbf{Problem 43}:\\

(a): 
$$
E(r) = \begin{cases}
	0 & r<b\\
	k(2q)/r^2 & b<r<c\\
	0 & c<r<d\\
	k(6q)/r^2 & d<r.
\end{cases}
$$
oriented outward from the origin everywhere it is nonzero.\\

(b): The charges at $r=a,b,c,d$ are $0,2q,-2q,6q$ respectively.\\

\textbf{Problem 44}:\\

(a): 
$$
E(r) = \begin{cases}
	0 & r<b\\
	k(2q)/r^2 & b<r<c\\
	0 & c<r
\end{cases}
$$
oriented outward from the origin everywhere it is nonzero.\\

(b): The charges at $r=a,b,c,d$ are $0,2q,-2q,0$ respectively.\\

\textbf{Problem 46}:\\

(a): Because the electric field inside the solid sphere nucleus points outward everywhere except where it is 0 at the center, so an electron (because it is negatively charged) moves in the opposite direction of the field, toward the center, where it is fixed.\\

(b): Simple harmonic motion occurs when the restoring force is proportional to the displacement from equilibrium. In this case, if the displacement is $r<R$, the electric field magnitude is $kQr/R^3$, and the force is $kr/R^3$ directed toward the center. This is proportional to $r$ with constant $k/R^3$ so it is SHM. The frequency of oscillation can be calculated by
$$
f = \frac{1}{2\pi} \sqrt{\frac{k_{spring}}{m}} = \frac{1}{2\pi} \sqrt{\frac{k}{R^3m}}.
$$

(c):
$$
4.57 \cdot 10^{14} = \frac{1}{2\pi} \sqrt{\frac{k}{R^3m}} \implies R = 9.4\cdot 10^{-4} m
$$
This is way larger than any atom's actual radius.\\

(d): If the displacement was greater than $R$, it would be pulled into the nucleus, but the motion would not be simple harmonic while it was outside.\\

\textbf{Problem 53}:\\

(a): If $d$ is the distance from the center of the sphere, then the electric field at this distance is 
$$
kQ/d^2 = \frac{k(\tfrac43 \pi d^3)\rho}{d^2} = k\tfrac43 \pi d \rho = \frac{d\rho}{3\ep_0}
$$
and $d=|\vec{r}-\vec{b}|$ so this is the answer expected.\\

(b): The electric field of the entire sphere of radius $R$ is $k\tfrac43 \pi \vec{r}\rho$, and the electric field of the little sphere of radius $a$ is $k\tfrac43 \pi (\vec{r}-\vec{b})\rho$ inside of it, so by superposition we can subtract to get the electric field of the sphere with the hole:
$$
k\tfrac43 \pi\vec{r}\rho - k\tfrac43 \pi(\vec{r}-\vec{b}) = k\tfrac43 \pi\vec{b}\rho 
$$
which does not depend on $\vec{r}$, thus it is uniform inside the hole.\\

\textbf{Problem 56}:\\

The electric field in an insulating sphere is $k\tfrac43 \pi r \rho$. So the slope wrt $r$ is $k\tfrac43 \pi \rho$, and in the graph this slope is $(3\cdot 10^4)/(4\cdot 10^{-3})$, so we have
$$
k\tfrac43 \pi \rho = \frac{3\cdot 10^4}{4\cdot 10^{-3}} \implies \rho = 2\cdot 10^{-4} C/m^3.
$$

\end{document}